\chapter{开篇辞}

\begin{quoting}怀悲悯与清凉之心，以慧灯驱除愚痴之暗冥，\\为俱有人与不死\footnote{不死 \textit{amara}：这里是指天人。}的世间所重，我礼拜解脱于趣的\textbf{善逝}。\end{quoting}

\begin{quoting}佛陀亦要修习并且证得的觉悟之相，\\趋近之便离垢，我礼拜这无上的\textbf{法}。\end{quoting}

\begin{quoting}善逝的孳乳子嗣，魔军的撼动者，\\且数众有八，我以头礼拜\textbf{圣僧伽}。\end{quoting}

\begin{quoting}如是，对净喜之觉的我，以这礼拜三宝所造福德\\之威力，既已善驱除了障碍。\end{quoting}

\begin{quoting}精微的长篇契经，\textbf{高贵的长阿含}，\\为佛与随佛所赞叹，有带来信的功德，\footnote{四部义注的「开篇辞」，唯此颂及末句不同。}\end{quoting}

\begin{quoting}旨在阐明其义的\textbf{义注}，在最初由五百\\自在\footnote{自在 \textit{vasin}：这里是指阿罗汉，下文「大摩哂陀自在」亦同。}结集，后来又再结集，\end{quoting}

\begin{quoting}然后，被大摩哂陀自在带至僧伽罗岛，\\存于僧伽罗语中，为了岛民们的义利。\end{quoting}

\begin{quoting}随后，我移除了僧伽罗语，将悦意的语言，\\与经典的理法相适、离去过失者引入，\end{quoting}

\begin{quoting}不歧异于上座传统之灯的诸长老、\\极精微裁断的大寺住者们\footnote{大寺住者们 \textit{Mahāvihāre nivāsīnaṃ}：PTS 本作「大寺认可者们 \textit{Mahāvihārādhivāsīnaṃ}」。}的教义，\end{quoting}

\begin{quoting}经舍弃重复出现之义，我将阐明义，\\为了善人的满足，还为了法的久住。\end{quoting}

\begin{quoting}戒论、头陀法，以及所有的业处，\\伴以性行的分类，禅那与等至的详说，\end{quoting}

\begin{quoting}所有的神通，以及慧的总体确定，\\蕴、界、处、根，以及四圣\end{quoting}

\begin{quoting}谛，缘的行相的开示，极遍净微妙的方法，\\未解脱的经典之道，以及毗婆舍那的修习，\end{quoting}

\begin{quoting}如是一切，因为在\textbf{清净道论}中已由我极遍净地\\讲述，所以，我在此将不再考察彼等。\end{quoting}

\begin{quoting}这\textbf{清净道论}既然也在四阿含之中\\存在，它将在那里阐明如上所说之义，\end{quoting}

\begin{quoting}而如是已作，所以，把握了彼等，再伴以这\\义注，便能了知依止长阿含之义。\end{quoting}

这里，\textbf{长阿含}者，从品上说，有戒蕴品、大品、波梨品等三品，从经上摄三十四经。在其品中，戒蕴品为初，经中则为梵网。而梵网之初的「如是我闻」，由阿难尊者在第一次大结集时所说，为因缘之初。

\chapter{论第一次大结集}

\subsection*{大迦叶尊者起意\footnote{本章中的小标题均为译者所加。}}

且此第一次大结集者，虽然在律藏中被引入经典\footnote{即\textbf{律藏}·小品·五百犍度。}，但为了善巧于因缘，在此亦应当知。当以转法轮为初，直至游行者须跋陀之调伏，尽完佛陀的义务，世尊、世间之依怙于拘尸那罗近旁、众末罗之娑罗林中的娑罗双树间，在毗舍佉满月日的黎明时分，以无余依涅槃界般涅槃后，在分配世尊的遗骨之日，聚集的七十万比丘中的僧团长老、大迦叶尊者忆念起在世尊般涅槃七天时出家耆宿须跋陀所说之语：

\begin{quoting}且止！朋友！莫忧！莫悲！我等善解脱于大沙门，且我等备受压迫——「这对你们许可、这对你们不许可」，但现在我们想做什么便做什么，想不做什么便不做什么。{\bluefootnotesize 〔小品第 437 段〕}\end{quoting}

并且认为僧团再难如此聚集，思惟道：「然而有这可能，即恶比丘认为『宣教是过去的大师的』，得了偏党，不久便使正法消逝，但只要法律尚存，宣教便非过去的大师的。因为如世尊所说：

\begin{quoting}阿难！我对尔等开示、施设的法与律，在我身后便是尔等的大师。{\bluefootnotesize 〔长部第 2.216 段〕}\end{quoting}

我何不结集法与律，好让这教法成为历时的久住者？且我在世尊说了

\begin{quoting}那你，迦叶！会持我破旧的麻布粪扫衣吗？{\bluefootnotesize 〔相应部第 2.154 段〕}\end{quoting}

之后，便被共同受用衣所摄受\footnote{所摄受 \textit{anuggahito}：据 PTS 本补，缅甸本与下句共用，亦通。}，并以

\begin{quoting}诸比丘！只要我希望，我便能离诸爱欲、离诸不善法，进入并住于有寻有伺、因远离而生喜乐的初禅，诸比丘！只要迦叶希望，他便能离诸爱欲、离诸不善法，进入并住于有寻有伺、因远离而生喜乐的初禅。{\bluefootnotesize 〔相应部第 2.152 段〕}\end{quoting}

等方法，被于九次第住、六神通等类的上人法置于与自己完全相同所摄受，同样，还因如空中挥掌的不固著之心和如月般的行道\footnote{二喻见\textbf{相应部}第 2.146 段。}受到赞赏，对此我\footnote{我 \textit{me}：据 PTS 本补。}责无旁贷！难道世尊不是认为『他将是延续正法世系者』，好似国王交付自身的铠甲和权力给延续自己家族世系的子嗣，以此不共的摄受摄受我，并以此显赫的赞赏赞赏我吗？」便为了结集法律，令众比丘生起勇猛。如说：

\begin{quoting}于是，大迦叶告诸比丘：「一时，朋友！我与大比丘僧团五百比丘一起，行于从波婆到拘尸那罗的旅途……」\footnote{此即\textbf{律藏}·小品·五百犍度的开篇。}{\bluefootnotesize 〔小品第 437 段〕}\end{quoting}

一切当知如「须跋陀章\footnote{须跋陀章 \textit{Subhadda-kaṇḍa}：现在的律藏和大般涅槃经中并无此标题，当即上述的「出家耆宿须跋陀」事。}」详说，而我们将在「大般涅槃」末尾的出现处谈论其义。

\subsection*{拣选比丘}

随后，他更说：

\begin{quoting}噫！朋友！让我们结集法与律！先前非法闪耀，法被压制，先前非律闪耀，律被压制，先前说非法者有力，说法者乏力，先前说非律者有力，说律者乏力。{\bluefootnotesize 〔小品第 437 段〕}\end{quoting}

比丘们便说：

\begin{quoting}那么，尊者！请长老拣选众比丘！{\bluefootnotesize 〔小品第 437 段〕}\end{quoting}

然而，长老排除了整个九分大师教的圣典持者，以及数百数千的凡夫、须陀洹、斯陀含、阿那含、干观者、漏尽比丘，只得到三藏一切圣典等的持者、证无碍解、有大威力、多数被世尊置于上首的三明等的漏尽比丘五百减一。就此而说：

\begin{quoting}于是，大迦叶尊者便拣选了五百减一个阿罗汉。{\bluefootnotesize 〔小品第 437 段〕}\end{quoting}

那为何长老要留一空缺呢？为了给阿难尊者创造机会。因为不与阿难尊者一起，便不能作法的结集，而此尊者是有学，还俱应作，所以不能一起。但因为没有任何十力开示的契经、应颂等是未经其现量的，如说：

\begin{quoting}我从佛陀取了八万二千，从比丘二千，\\由我转起之法，计八万四千。{\bluefootnotesize 〔长老偈第 1027 颂〕}\end{quoting}

所以也不能去除。

如果即便作为有学，由对法的结集多有助益故，应被长老拣选，那又为什么未被拣选呢？由避免他人的指责故。因为长老极度信赖阿难尊者，同样，即便头上生白，仍以

\begin{quoting}此童子尚不知量。{\bluefootnotesize 〔相应部第 2.154 段〕}\end{quoting}

之童子论来教诫他，且尊者出生自释迦家族，是如来之弟、叔父之子。对此，某些比丘认为似是欲的邪向，便指责道：「搁置了众多证得无学无碍解的比丘，长老拣选了证得有学无碍解的阿难。」为避免这他人的指责，便以「去除阿难，不能作法的结集，我将以众比丘的许可来获取他」，未作拣选。

于是，众比丘自己便为了阿难请求长老，如说：

\begin{quoting}众比丘便对大迦叶尊者说：「尊者！这阿难尊者虽然是有学，不会因欲、嗔、痴、怖畏至于非趣，并以此在世尊跟前学习了众多的法与律，因此，尊者！请长老也拣选阿难尊者！」于是，大迦叶尊者便也拣选了阿难尊者。{\bluefootnotesize 〔小品第 437 段〕}\end{quoting}

如是，以众比丘的许可，与被拣选的这尊者一起，便成五百长老。

\subsection*{大迦叶尊者作羯磨语}

\begin{quoting}于是，众长老比丘便想：「我们能到哪里去结集法与律呢？」然后，众长老比丘想道：「王舍城行处甚大、坐卧处甚多，我们何不去在王舍城住而安居、结集法与律呢？其他比丘则不应在王舍城度安居。」{\bluefootnotesize 〔小品第 437 段〕}\end{quoting}

但他们为什么这样想？「这是我们的持久羯磨，任何异分补特伽罗进入僧团之中，将会打扰\footnote{打扰 \textit{ukkoṭeti}：亦有「重开已决之事」的意思。}。」

于是，大迦叶尊者便以白二羯磨宣告：

\begin{quoting}请僧团听我！朋友！如果僧团准备完毕，僧团应共许这五百比丘在王舍城住而安居，以便结集法与律，不应与其他比丘在王舍城住而安居。\end{quoting}

\begin{quoting}这是告白。\end{quoting}

\begin{quoting}请僧团听我！朋友！共许这五百比丘在王舍城住而安居，以便结集法与律，不应与其他比丘在王舍城住而安居。\end{quoting}

\begin{quoting}若尊者认可所共许的这五百比丘「在王舍城住而安居，以便结集法与律，不应与其他比丘在王舍城住而安居」，他应默然，若不认可，他应说。\end{quoting}

\begin{quoting}为僧团所共许的这五百比丘「在王舍城住而安居，以便结集法与律，不应与其他比丘在王舍城住而安居」，僧团已认可，所以默然，我如是持之。{\bluefootnotesize 〔小品第 438 段〕}\end{quoting}

而这羯磨语是在如来般涅槃后的第二十一天所作。因为世尊在毗舍佉的满月日的黎明时分般涅槃，于是此后的七天，人们以香鬘等供养金色之身，如是七天便得名善庆祝日，随后的七天，以火葬柴堆之火荼毗，以七天造了矛垣，在议事厅\footnote{议事厅 \textit{santhāgāra}：原文作 sandhāgāra，兹从 PTS 本。}的会堂供养了遗骨，便过了二十一天，在逝瑟吒月\footnote{逝瑟吒月 \textit{Jeṭṭhamūla}：即毗舍佉之后的一月，公历五～六月。}白分的第五日，分配了遗骨。在此分配遗骨之日，向聚集的大比丘僧团告知了出家耆宿须跋陀所作的非正行，并在以所述之法拣选了众比丘后，作了这羯磨语。

且作了这羯磨语后，长老告诸比丘：「朋友！现在，你们有四十天的空闲，此后便不得说『我们有这般障碍』，所以在这期间，若有疾病的障碍，或阿阇黎、亲教师的障碍，或父母的障碍，或钵应烤，或衣应作，他应断了这障碍，作此应作！」

且如是说罢，长老为自己的五百随从围绕，去了王舍城。其他的大长老们也带了各自的眷属，欲安慰罹患忧箭的大众，前往各自的方向。而富楼那长老则为七百比丘围绕，想「我当安慰陆续来到如来般涅槃处的大众」，仍留在拘尸那罗。

\subsection*{阿难尊者行略}

阿难尊者仍如先前未入灭般，亲自持了已入灭的世尊的衣钵，与五百比丘一起，前往舍卫国游行。而随着前行，其随从的比丘不可胜计。在这尊者所到之处，即起大悲泣。而当长老渐次到达舍卫国时，舍卫国的居民听闻「据说长老来了」后，手持香鬘等迎接，说着「尊者！阿难！先前您与世尊一起前来，今天把世尊置于何处而来」，便嚎啕大哭。即如佛世尊般涅槃的当天一般，起大悲泣。

对此，阿难尊者以与无常等相应的法论抚慰大众后，进入祇园，礼拜了十力曾住的香房，打开门，取出床椅拍打，扫除香房，丢掉凋萎的花鬘和垃圾，再放回床椅，置于原处，作了世尊住时一切应作的义务。且当作时，在扫除浴槽、供水等时，礼拜了香房，便以「世尊！这是您沐浴之时，这是开示法之时，这是给予比丘们教诫之时，这是作狮子卧之时，这是洗脸之时」等方式而作悲泣，正如以对世尊功德聚的甘露味之知所建立的情感，又如非漏尽者，又如在数百千生中对彼此的资助所生的心的温柔。

某位天人便警示他道：「尊者！阿难！你如是悲泣，如何能安慰其他人呢？」他因他的话语而心生警惕，克制后，为恢复从如来般涅槃起由久站久坐之故界已增盛的身体，在第二天喝下乳泻药后，唯坐于寺中。就此，便对须婆学童所遣的学童说：

\begin{quoting}非时，学童！我今天只喝了药，也许明天我们会前往。{\bluefootnotesize 〔长部第 1.447 段〕}\end{quoting}

第二天，以支多迦长老为随从沙门，一起前去，就须婆学童之所问，说了这长部中名为\textbf{须婆经}的第十经。

\subsection*{五百比丘入安居}

于是，阿难长老教人修葺了祇园大寺中的破损残败，在接近雨安居开始时，舍下比丘僧团，到达王舍城，而其他结集法的比丘们也同样。如是到达时，就他们而说：

\begin{quoting}于是，长老比丘们便到了王舍城，为结集法与律。{\bluefootnotesize 〔小品第 438 段〕}\end{quoting}

他们作了阿沙陀\footnote{阿沙陀 \textit{Āsāḷhī}：即逝瑟吒之后的一月，公历六～七月。}满月日的布萨，在半月的第一天聚集后，便入了安居。

而那时，有十八座大寺围绕着王舍城，它们全都荒废、败落、尘封。因为当世尊般涅槃，所有比丘都持了各自的衣钵，丢弃寺庙与僧寮而去。长老们为了供养世尊的话语，并为了摆脱外道的论议，在此订立规约，便思惟：「让我们在第一个月修葺破损残败！」因为外道们会如是说：「沙门乔达摩的弟子们在大师尚在时照料寺庙，在般涅槃时则予丢弃，许多家族遍舍的巨财都亡失了。」即是说为了摆脱他们的论议作此思惟。且经如是思惟，便订立了规约。就此而说：

\begin{quoting}于是，长老比丘们想道：「朋友！修葺破损残败为世尊所赞叹，噫！朋友！让我们在第一个月修葺破损残败，在中间一月聚集后，我们将结集法与律！」{\bluefootnotesize 〔小品第 438 段〕}\end{quoting}

他们在第二天去到国王前而立。国王前来礼拜，便问自己当作的义务：「尊者！你们为何前来？」长老们为了修葺十八大寺，报告了劳务。国王便给予劳务的人工。长老们在第一个月教人修葺完所有寺庙，便告知国王：「大王！修葺已毕，现在我们要结集法律！」「善哉！尊者！尽管去做！让我有敕命之轮，而你们有法轮！尊者！命令我该做什么？」「行结集比丘们的共坐之处，大王！」「我在哪里做呢？尊者！」「在吠婆罗\footnote{吠婆罗 \textit{Vebhāra}：即环绕王舍城的五山之一。}山间的七叶窟前适合去做，大王！」

「善哉！尊者！」阿阇世王教人建了圆亭，仿佛工艺天所化现，墙柱阶梯善加规划，装饰以多种鬘饰、藤饰，似超迈了王家居处的盛况，似讪笑着天宫的辉光，似吉祥的居所，又似人天目光之飞鸟齐栖的津渡\footnote{即人天瞩目之意。}，似世间可意之物凝聚的可观的醍醐精华，庄严以悬挂各种花环并飘谢着的怡人华盖，彩色的地坪以多种花夯实，善加竣工，似以众宝和彩色摩尼镶嵌的地面，恍若梵宫，在此大圆亭中，为五百比丘设了五百张无价但经许可的铺盖，依于南侧，面朝北方，为长老们而设，在圆亭中间，面朝东方，为佛世尊设了应坐的法座，并在此放置了嵌有象牙的扇子，便教人告知比丘僧团：「尊者！我的义务已毕。」

\subsection*{阿难尊者证阿罗汉}

而在那天，有些比丘就阿难尊者作如是说：「在此比丘僧团中，一位比丘散发着臭味而行。」长老听后，想「在此比丘僧团中，并无其他散发着臭味而行的比丘，他们肯定是在说我吧」，便生起悚惧。有些人还对他说：「朋友！阿难！明天集会，而你还是有学，还俱应作，因此你去集会不合适，请勿放逸！」

于是，阿难尊者想「明天集会，我要是作为有学去集会不合适」，彻夜以身至念度过，在夜里的黎明时分从经行道下来，进入寺庙，便以「我将卧倒」转向于身体，双足从地抬起，而头未碰枕头，心便在此之间以无取从诸漏解脱。

因为这尊者在外以经行度夜，不能转起殊胜，便思惟：「世尊不是对我说『阿难！你已积累福德，应专心于至上，速速成就无漏』{\bluefootnotesize 〔长部第 2.207 段〕}吗？而诸佛的谈论无有过失，我的精进太过，因此我的心导致掉举，噫！我要驱使精进平衡！」便从经行道下来，站在洗脚处洗了脚，进入寺庙，坐在床上，以「休息片刻」向床上放下身体。双足从地抬起，而头未碰枕头，心在此之间以无取从诸漏解脱，长老的阿罗汉离四威仪。因此，当说「此教法中，哪位比丘非卧、非坐、非站、非经行证得阿罗汉」时，说「阿难长老」是合适的。

于是，长老比丘们在第二天，黑分的第五天，食事已毕，收纳了衣钵，便聚集在法堂。然后，阿难尊者作为阿罗汉，便前去集会。他如何前去？他想「现在，我值得进入集会」，心里高兴、满足，把衣偏覆一肩，如脱茎的棕榈果，如掉在黄毯上的天然摩尼，如拨开云翳露出的满月，如初阳照射、花粉飘散、胎室赤黄的莲花，并以遍净洁白、具有光明祥瑞的高贵面容告知着自己证得阿罗汉般前去。

于是，大迦叶尊者见到后便想：「证得阿罗汉的阿难君真是光彩！如果大师尚在，今天肯定会为阿难鼓掌，噫！现在，我来给予应由大师给的掌声吧！」便鼓掌三次。

而中部诵者们却说：「阿难长老欲通知自己证得阿罗汉，便未与比丘们同往，比丘们按长幼坐到各自的坐处，空着阿难长老的坐处而坐。对此，有人作如是说：『这坐处是谁的？』『阿难的。』『那阿难去哪里了？』此时，长老便思惟：『现在是我前来的时刻。』随后，为显示自己的威力，潜入地下，刚好在自己的坐处现身。」也有人说是从空中前来而坐。无论如何，总之，大迦叶尊者见到后给予掌声是恰当的。

\subsection*{结集律}

而当他如是到来时，尊者大迦叶长老便告诸比丘：「朋友！我们先结集什么？法还是律呢？」比丘们便说：「尊者！大迦叶！\textbf{律者}，佛教之寿命也，律存则教法方存，所以，让我们先结集律！」「谁担此任？」「优婆离尊者。」「难道阿难不能？」「非不能！而是正等正觉者尚在时，曾就律的学习置优婆离尊者于上首：

\begin{quoting}诸比丘！我的声闻比丘中，持律之上首，即优婆离是。{\bluefootnotesize 〔增支部第 1.228 段〕}\end{quoting}

所以，让我们询问优婆离长老，结集律！」随后，长老便自己同意\footnote{同意 \textit{sammanni}：或译作「荐举、选择」，后二处同。}自己去询问律，优婆离长老也同意去作解答。

此处的圣典为：

\begin{quoting}于是，大迦叶尊者通知僧团：「请僧团听我！朋友！如果僧团准备完毕，我将询问优婆离律。」\end{quoting}

优婆离尊者也通知僧团：

\begin{quoting}请僧团听我！尊者！如果僧团准备完毕，我将解答由大迦叶尊者询问的律。{\bluefootnotesize 〔小品第 439 段〕}\end{quoting}

如是同意自己后，优婆离尊者从坐起，把衣偏覆一肩，礼拜了长老比丘们，便坐于法座，拿起嵌有象牙的扇子。随后，大迦叶长老坐于长老之座，便询问优婆离尊者律：

\begin{quoting}「朋友优婆离！第一波罗夷在何处制？」\\「在毗舍离，尊者！」\\「关于谁？」\\「关于须提那·迦兰陀子。」\\「于何事由？」\\「于淫欲法。」\\然后，大迦叶尊者便询问第一波罗夷的事由，询问因缘，询问人，询问制，询问随制，询问犯，询问无犯。{\bluefootnotesize 〔小品第 439 段〕}\end{quoting}

优婆离尊者便随问而答。

那么，在第一波罗夷中是否有任何应当删除或补充的吗？没有应当删除的。因为佛世尊之所说并无应当删除的。诸如来甚至不说一字无意义者。但声闻或诸天之所说则有应当删除的，结集法的长老们便删除之。而应当补充的则所在皆有，所以，在适于补充之处，他们便予补充。那有哪些？若「尔时」、若「而在尔时」、若「于是」、若「如是说罢」、若「他便说」，如是等仅作连接语者。如是补充了适于补充处之后，他们便确立了「此即第一波罗夷」。

当第一波罗夷被引入结集后，五百阿罗汉便按引入结集之法合诵：「尔时，佛世尊住于毗兰若……」当他们刚开始诵时，似有掌声鼓起，大地以水为边界而震动。

仍按此法，他们将其余的三波罗夷引入结集，确立了「此即波罗夷章」。确立了

\begin{itemize}[nosep]
    \item 十三僧残为「十三法」，
    \item 二学处为「不定」，
    \item 三十学处为「舍波夷提」，
    \item 九十二学处为「波夷提」，
    \item 四学处为「悔过」，
    \item 七十五学处为「众学」，
    \item 七法为「灭诤」。
\end{itemize}

如是宣告确立了二百二十七学处为「大分别」。在大分别的终了，大地仍按前述之法而震动。

随后，在比丘尼分别中，他们确立了

\begin{itemize}[nosep]
    \item 八学处为「此即波罗夷章」，
    \item 十七学处为「十七法」，
    \item 三十学处为「舍波夷提」，
    \item 一百六十六学处为「波夷提」，
    \item 八学处为「悔过」，
    \item 七十五学处为「众学」，
    \item 七法为「灭诤」。
\end{itemize}

如是宣告三百又四学处为「比丘尼分别」，确立了「此二部分别计六十四诵\footnote{诵 \textit{bhāṇavāra}：文本长度的单位，类似于我国的「卷」，巴利语法「\textbf{声律} \textit{Saddanīti}」第 358 页说：「对以三十二音节构成者，二百五十成为一诵，实八千音节。」}」。在二部分别的终了，按前述之法，亦有大地震动。

仍按此方法，将八十诵之量的犍度、二十五诵之量的附随引入结集，确立了「此即\textbf{律藏}\footnote{如此，则律藏共计 169 诵。}」。在律藏的终了，按前述之法，亦有大地震动。

他们便将其托付给优婆离尊者：「朋友！请将此传授给你的依止者！」在律藏结集的终了，优婆离长老放下嵌有象牙的扇子，从法座下来，礼拜了长老比丘们，便坐于自己的坐处。

\subsection*{结集法}

结集完律后，大迦叶尊者欲结集法，便问比丘们：「当结集法时，应令谁人担此任来结集法？」比丘们便说：「应令阿难长老担此任。」

于是，大迦叶尊者通知僧团：

\begin{quoting}请僧团听我！朋友！如果僧团准备完毕，我将询问阿难法。\end{quoting}

阿难尊者通知僧团：

\begin{quoting}请僧团听我！尊者！如果僧团准备完毕，我将解答由大迦叶尊者询问的法。\end{quoting}

于是，阿难尊者从坐起，把衣偏覆一肩，礼拜了长老比丘们，便坐于法座，拿起嵌有象牙的扇子。然后，大迦叶长老便问比丘们：「朋友！我们先结集哪一藏？」「经藏，尊者！」「经藏中有四等诵\footnote{等诵 \textit{saṅgīti}：也有「结集」之义，这里应指经藏的四部分。}，先结集其中哪一等诵？」「长等诵，尊者！」「长等诵中有三十四经、三品，先结集其中哪一品？」「戒蕴品，尊者！」「戒蕴品中有十三经，先结集其中哪一经？」「尊者！梵网经者，庄严以三种戒，摧破种种邪命、诡诈、虚谈等，解开六十二见之网，震动一万世界，让我们先结集它！」

于是，大迦叶尊者便对阿难尊者说：

\begin{quoting}「朋友阿难！梵网是在何处所说？」\\「尊者！在王舍城与那烂陀之间的王家芒果树园。」\\「关于谁？」\\「游行者善念\footnote{善念 \textit{Suppiya}：译名从\textbf{长阿含经}·梵动经。}和学童梵赐。」\\「于何事由？」\\「于赞誉与毁谤。」\\然后，大迦叶尊者便询问阿难尊者梵网的因缘，询问人，询问事由。{\bluefootnotesize 〔小品第 440 段〕}\end{quoting}

阿难尊者便作解答。在解答的终了，五百阿罗汉便合诵。且仍按前述之法，有地震动。

如是结集了梵网，随后更以「那么，朋友阿难！沙门果是在何处所说」等方法，以次第的问答，连同梵网一起，结集了所有十三部经，宣告确立了「此即戒蕴品」。

此后是大品，此后是波梨品，如是结集了三品所摄、三十四经组成、六十四诵之量的经典，说「此即\textbf{长部}」，便托付给阿难尊者：「朋友！请将此传授给你的依止者！」

此后，结集了八十诵之量的\textbf{中部}，便托付给法将舍利弗长老的依止者：「请你们保存它！」

此后，结集了一百诵之量的\textbf{相应部}，便托付给大迦叶长老：「尊者！请将此传授给您的依止者！」

此后，结集了二十百诵之量的\textbf{增支部}，便托付阿那律长老：「请将此传授给您的依止者！」

\subsection*{结集阿毗达摩}

此后，结集了被称为「法集、分别、界论、人施设、论事、双、发趣」的阿毗达摩，被赞赏为精微之智领域的经典，说「此即\textbf{阿毗达摩藏}」，五百阿罗汉便合诵。仍按前述之法，有地震动。

随后，又结集了本生、义释、无碍解道、譬喻、经集、小诵、法句、自说、如是语、天宫事、饿鬼事、长老偈、长老尼偈等经典，说「此即\textbf{杂典}\footnote{杂典：即现在的\textbf{小部}，这里的原文作 Khuddaka-gantha，Khuddaka 有小、杂的意思，gantha 则是典籍之义。}」，长部诵者们说：「它们也被引入结集于阿毗达摩藏中。」而中部诵者们则说：「连同所行藏、佛种姓一起，这一切都名为杂典，系属于经藏。」

\subsection*{佛语之分类}

如是，这一切佛语依味为一种，依法律为二种，依初中后为三种，依藏也同样，依部为五种，依分为九种，依法蕴为八万四千种。

如何\textbf{依味为一种}？因为它由世尊于无上正等正觉现等觉已，直至以无余依涅槃界般涅槃，其间四十五年对天、人、龙、夜叉等以教授或以省察而说，这所有唯解脱味一味。如是依味为一种。

如何\textbf{依法律为二种}？且这一切都归于法与律。这里，律藏即律，其余的佛语为法。正是因此而说「我们何不结集法与律」，以及「我要问优婆离律，问阿难法」。如是依法律为二种。

如何\textbf{依初中后为三种}？因为这一切都以最初的佛语、中间的佛语、最后的佛语为三种。这里，

\begin{quoting}我曾不停地流转于多生的轮回，\\寻觅着造屋者，苦哉再再出生！\end{quoting}

\begin{quoting}造屋者！你已被发现，将不再造屋！\\你的椽桷都已断裂，屋脊也被拆除，\\心已到达离行，证得了渴爱的灭尽。{\bluefootnotesize 〔法句第 153-54 颂〕}\end{quoting}

这即最初的佛语。有些人说是犍度中的自说之颂，即「当法确乎显现之时……」{\bluefootnotesize 〔大品第 1 段〕}。然而，当知这是在半月的第一天，证得一切知性者以喜悦所造之智省察缘的行相生发的自说之颂。

而在般涅槃时，他说：

\begin{quoting}噫！现在，诸比丘！我昭告你们：「诸行是坏灭法，尔等应以不放逸而成就！」{\bluefootnotesize 〔长部第 2.218 段〕}\end{quoting}

这即最后的佛语。二者之间所说的，名为中间的佛语。如是，依最初、中间、最后的佛语为三种。

如何\textbf{依藏为三种}？这一切还以律藏、经藏、阿毗达摩藏为三种。这里，汇集了第一次结集中结集与未结集的一切，二波罗提木叉、二分别、二十二犍度、十六附随等，名为\textbf{律藏}。梵网等三十四经所摄为长部，根本方法经等一百五十经所摄为中部，度暴流经等七千七百六十二经所摄为相应部，心之占据经等九千五百五十七经所摄为增支部，依小诵、法句、自说、如是语、经集、天宫事、饿鬼事、长老偈、长老尼偈、本生、义释、无碍解道、譬喻、佛种姓、所行藏的十五种为小部，名为\textbf{经藏}。法集、分别、界论、人施设、论事、双、发趣，名为\textbf{阿毗达摩藏}。

\subsection*{释律字}

这里，

\begin{quoting}由种种和特殊的方法故，且由调伏身语故，\\此调伏被明了调伏之义者称为律。\footnote{这是从语源解释「律 \textit{vinaya}」：取了「种种 \textit{vividha}」和「特殊 \textit{visesa}」中的 vi，与「方法 \textit{naya}」组合成 vinaya。}\end{quoting}

此中的种种，即五分波罗提木叉诵、波罗夷等七聚罪、总纲和分别等类的方法。且有作为特殊的，旨在令其从严和从宽的随制的方法。并由禁止身语的过行故，调伏了身语。所以，由种种方法故、特殊方法故、调伏身语故，被称为律。因此，这是为了善巧其语义而说：

\begin{quoting}由种种和特殊的方法故，且由调伏身语故，\\此调伏被明了调伏之义者称为律。\end{quoting}

\subsection*{释经字}

而另一为：

\begin{quoting}由对诸义利指示故、善说故、产生故、流淌故、\\善庇护故，且与线同分故，被称为经。\footnote{这也是从语源解释「经 \textit{sutta}」。}\end{quoting}

因为它指示自身的义利、他人的义利等类的义利。且此中的义利为善说，由随顺被调伏者的意乐而说故。且它产生义利，如谷物产出果实一般。且它流淌，如乳牛令乳滴淌一般。且它善加庇护、守护彼等。且它与线同分，因为好比线对于木匠是度量一般，如是，它对于智者也是。又好比被线贯串的花不会被零落、不会被拆散，如是，被它贯串的义利也是。因此，这是为了善巧其语义而说：

\begin{quoting}由对诸义利指示故、善说故、产生故、流淌故、\\善庇护故，且与线同分故，被称为经。\end{quoting}

\subsection*{释阿毗达摩字}

而另一为：

\begin{quoting}凡此中具有增长、俱相、受尊敬、明确、\\以及所说的超越的法，因此被称为阿毗达摩。\end{quoting}

因为这「阿毗 {\footnotesize \textit{abhi}}」一字见于增长、相、受尊敬、明确、超越等处。如在

\begin{quoting}我有剧烈的苦受在增长 {\footnotesize \textit{abhikkamanti}}，而非消退。{\bluefootnotesize 〔中部第 3.389 段〕}\end{quoting}

等处作增长，在

\begin{quoting}凡那些知名而特殊 {\footnotesize \textit{abhiññātā abhilakkhitā}} 的夜晚。{\bluefootnotesize 〔中部第 1.49 段〕}\end{quoting}

等处作俱相，在

\begin{quoting}王中胜王 {\footnotesize \textit{abhirājā}}，人中因陀。{\bluefootnotesize 〔中部第 2.399 段〕}\end{quoting}

等处作受尊敬，在

\begin{quoting}堪能调伏于增上法、增上律 {\footnotesize \textit{abhidhamme abhivinaye}}。{\bluefootnotesize 〔大品第 85 段〕}\end{quoting}

等处作明确——即是说于彼此不混杂的法与律中，在

\begin{quoting}以殊胜 {\footnotesize \textit{abhikkantena}} 的容貌。{\bluefootnotesize 〔天宫事第 819 颂〕}\end{quoting}

等处作超越。\footnote{以上是说明「阿毗」一字之义，下面则举阿毗达摩藏中的例子来说明「具有增长、俱相、受尊敬、明确、超越之法」。}且此中，以

\begin{quoting}为投生于色而修习道。{\bluefootnotesize 〔法集论第 251 段〕}\\以与慈俱行之心遍满一方而住。{\bluefootnotesize 〔分别论第 642 段〕}\end{quoting}

等方法而说为具有增长的法，以

\begin{quoting}色所缘，或声所缘……{\bluefootnotesize 〔法集论第 1 段〕}\end{quoting}

等方法，由所缘等可作相故，而说为俱相，以

\begin{quoting}有学法、无学法、出世间法。{\bluefootnotesize 〔法集论·总纲〕}\end{quoting}

等方法而说为受尊敬——即应得尊敬之意，以

\begin{quoting}有触，有受……{\bluefootnotesize 〔法集论第 1 段〕}\end{quoting}

等方法，由自性明确故，而说为明确，以

\begin{quoting}大界法、无量法……无上法。{\bluefootnotesize 〔法集论·总纲〕}\end{quoting}

等方法而说为超越的法。因此，这是为了善巧其语义而说：

\begin{quoting}凡此中具有增长、俱相、受尊敬、明确、\\以及所说的超越的法，因此被称为阿毗达摩。\end{quoting}

\subsection*{释藏字}

而此中未作区分者，其为：

\begin{quoting}藏，即明了藏义者从圣典及容器之义而说，\\与此组合后，应知有律等三。\end{quoting}

因为圣典在

\begin{quoting}莫因据有藏。{\bluefootnotesize 〔增支部第 3.66 段〕}\end{quoting}

等处也被称为「藏」，在

\begin{quoting}于是，若人取了锹、箧 {\footnotesize \textit{piṭaka}} 而来……{\bluefootnotesize 〔增支部第 3.70 段〕}\end{quoting}

等处亦为任何容器，所以「藏，即明了藏义者从圣典及容器之义而说」。现在，「与此组合后，应知有律等三」，如是，与此二种义的「藏」字复合后，律和此藏由圣典之相，以及由各自之义的容器，为律藏。仍按所述之法，经和此藏为经藏，阿毗达摩与此藏为阿毗达摩藏。如是，这些应知有律等三。

\subsection*{杂论三藏}

且如是了知已，又为了善巧此三藏中的种种品类：

\begin{quoting}其中如其所应，有开示、教法、论述等类，\\且能阐明修学、舍弃、甚深之相。\end{quoting}

\begin{quoting}学习的分类，以及在何处成与败，\\比丘如何圆满，亦能分别这一切。\end{quoting}

此处，有阐明，有分别。因为这三藏依次被称为敕命、世俗、第一义之开示，随犯、随顺、随法之教法，律仪与非律仪、见之解除、名色差别之论述。

此中，律藏，由有权敕命之世尊多从敕命开示故，为\textbf{敕命之开示}。经藏，由善巧于世俗之世尊多从世俗开示故，为\textbf{世俗之开示}。阿毗达摩藏，由善巧于第一义之世尊多从第一义开示故，为\textbf{第一义之开示}。

同样，第一，凡是多有违犯的有情，他们在此中随犯受教，为\textbf{随犯之教法}。第二，多种意乐、倾向、性行、信解的有情，在此中随顺受教，为\textbf{随顺之教法}。第三，只于法的积聚作「我、我所」之想的有情，在此中随法受教，为\textbf{随法之教法}。

同样，第一，此中作为对治过行而被论述的律仪与非律仪，为\textbf{律仪与非律仪之论述}——律仪与非律仪，即小与大的律仪，如业与非业、果与非果。第二，此中作为对治六十二见而被论述的见之解除，为\textbf{见之解除之论述}。第三，此中作为对治贪等而被论述的名色差别，为\textbf{名色差别之论述}。

且在此三者中，亦当知有三学、三舍弃与四种甚深之相。因为同样，律藏中特别说的是\textbf{增上戒学}，经藏中则是\textbf{增上心学}，阿毗达摩藏中则是\textbf{增上慧学}。

且律藏中为\textbf{违越的舍弃}，由对于烦恼，违越的对治属于戒故。经藏中为\textbf{缠的舍弃}，由缠的对治属于定故。阿毗达摩藏中为\textbf{随眠的舍弃}，由随眠的对治属于慧故。且第一为\textbf{彼分断}，另二为\textbf{镇伏、正断断}。第一为\textbf{舍弃恶行之杂染}，另二为\textbf{舍弃爱、见之杂染}。

且此中一一当知各有四种法、义、开示、通达甚深之相。这里，\textbf{法}，即经典。\textbf{义}，即其义。\textbf{开示}，即以心意为其确立的经典的开示。\textbf{通达}，即对经典及经典之义的如实理解。在此三者中都有此等法、义、开示、通达。因为好似大海之于兔等，对钝觉者难以跃过、不能立足，所以为甚深。如是，当知此中一一各有四种甚深之相。

另一方法为：\textbf{法}，即因，因为说「对于因之智为法无碍解」。\textbf{义}，即因之果，因为说「对于因之果之智为义无碍解」{\bluefootnotesize 〔分别论第 720 段〕}。\textbf{开示}，即施设，意即按照法来表达法，或以顺逆详略等论述。\textbf{通达}，即现观，且其从境域及不痴迷而有世间、出世间，随适义而理解诸法，随适法而理解诸义，随适施设之路而理解施设，或者，即彼彼随处所说诸法之应通达的被称为自相的无颠倒自性。

现在，因为在这些藏中的种种法，或种种义，以及照亮这无论如何应令知晓、现于听者之智的义的开示，与此中被称为无颠倒理解的通达，或者，即彼彼诸法之应通达的被称为自相的无颠倒自性，这一切对于未积累资粮的恶慧者，好似大海之于兔等，难以跃过、不能立足，所以为甚深。如是，亦当知此中一一各有四种甚深之相。

至此，

\begin{quoting}其中如其所应，有开示、教法、论述等类，\\且能阐明修学、舍弃、甚深之相。\end{quoting}

此颂之义已述。

\begin{quoting}学习的分类，以及在何处成与败，\\比丘如何圆满，亦能分别这一切。\end{quoting}

而此中，在三藏中，有三种学习的分类可见，即三种学习：蛇喻、为了出离、司库之学习。这里，凡错误把握，以诘难等因而掌握者，即\textbf{蛇喻}。就此而说：

\begin{quoting}譬如，诸比丘！有人希求蛇、寻觅蛇，遍求蛇而行，他若看见大蛇，在盘起的身体或尾巴处抓住它，则那蛇转回来咬他的手、或臂、或肢体的某一处，他便由此因遭受死亡，或死量之苦。这是什么原因？诸比丘！由错误把握蛇故。\end{quoting}

\begin{quoting}如是，诸比丘！于此，有些愚人掌握了法——契经……毗陀罗，他们掌握了这法后，不以慧考察这些法之义，他们便未认可这些不以慧考察的法之义，他们出于诘难的利益，以及出于投入讨论的利益掌握了法，且为此义利掌握法者，不能领受其义利，对于他们，这些法被错误把握，长夜转起非利与苦。这是什么原因？诸比丘！由错误把握诸法故。{\bluefootnotesize 〔中部第 1.238 段〕}\end{quoting}

然而，凡正确把握，只以希望戒蕴等的圆满而掌握者，非以诘难等因，即\textbf{为了出离}。就此而说：

\begin{quoting}对于他们，这些法被正确把握，长夜转起利益与乐。这是什么原因？诸比丘！由正确把握诸法故。{\bluefootnotesize 〔中部第 1.239 段〕}\end{quoting}

然而，凡已遍知蕴、已舍弃烦恼、已修习道、已通达不动、已证得灭的漏尽者，只是为了保护谱系、为了守护世系而掌握者，即\textbf{司库之学习}。

而比丘于律善行道者，依于戒的成就，圆满三明，且对这些在彼处作了分析。于经善行道者，依于定的成就，圆满六神通，且对这些在彼处作了分析。于阿毗达摩善行道者，依于慧的成就，圆满四无碍解，且对这些在彼处作了分析。如是，于这些善行道者，依次圆满这三明、六神通、四无碍解等类的成就。

而于律恶行道者，由与被许可的乐触的铺盖、披覆等触平等故，于被拒斥的有执受之触等处作无过想。即说：

\begin{quoting}我如此理解世尊开示之法，即凡这些障碍法，世尊所说之障碍，对受用彼等者，不足为障碍。{\bluefootnotesize 〔中部第 1.234 段〕}\end{quoting}

随后便圆满恶戒之相。于经恶行道者，不知晓在「诸比丘！存在、现存这四种补特伽罗……」{\bluefootnotesize 〔增支部第 4.5 段〕}等处之意趣而错误把握。就此而说：

\begin{quoting}以自己的错误把握，既诽谤我们，又伤害自己，并产生许多非福。{\bluefootnotesize 〔中部第 1.236 段〕}\end{quoting}

随后便圆满邪见性。于阿毗达摩恶行道者，过度驰骋于法思，且思惟诸多不应思惟者，随后便圆满心的散乱。即说：

\begin{quoting}诸比丘！这四种不应思惟者、不当思惟者，若思惟之，则疯狂、困扰。{\bluefootnotesize 〔增支部第 4.77 段〕}\end{quoting}

如是，于这些恶行道者，依次圆满这恶戒之相、邪见性、心的散乱等类的败坏。

至此，

\begin{quoting}学习的分类，以及在何处成与败，\\比丘如何圆满，亦能分别这一切。\end{quoting}

此颂之义已述。如是从种种品类了知了三藏，依于彼等，当知这佛语有三种。

\subsection*{五部}

如何\textbf{依部为五种}？且这一切以长部、中部、相应部、增支部、小部为五类。这里，何为\textbf{长部}？即三品所摄的梵网等三十四经。

\begin{quoting}三十四经，三品所摄，\\斯为长部，顺序之初。\end{quoting}

那为什么它被称为长部？对量长的经，由集合与住处故。因为集合、住处被称为部：

\begin{quoting}诸比丘！我不见另一部有如是之心，此即是，诸比丘！畜生生类。{\bluefootnotesize 〔相应部第 2.100 段〕}\end{quoting}

且此中，如「倾向之部、泥沼之部\footnote{倾向之部、泥沼之部 \textit{poṇikanikāyo cikkhallikanikāyo}：二词在三藏及注疏中仅见于解释「部」字。\textbf{疏}云「倾向、泥沼为刹帝利，他们的住处为倾向之部、泥沼之部 \textit{Poṇikā cikkhallikā ca khattiyā, tesaṃ nivāso poṇikanikāyo cikkhallikanikāyo ca}」，义有未安，下云「可作教法与世间上的例证」，则倾向之部当为教法的例证，泥沼之部为世间的例证。}」等可作教法与世间上的例证。其余部中的语义当知亦如是。

何为\textbf{中部}？即量中的五十品所摄的根本方法经等一百五十又二经。

\begin{quoting}彼处有一百五十经又二经，这\\中部为五十品所概括。\end{quoting}

何为\textbf{相应部}？即以天相应等所述的度暴流等七千七百六十二经。

\begin{quoting}七千经与七百经，\\并六十二经，此相应所摄。\end{quoting}

何为\textbf{增支部}？即以一一支增长所述的心之占据等九千五百五十七经。

\begin{quoting}九千经与五百经、\\五十七经，此即增支中数。\end{quoting}

何为\textbf{小部}？即整个律藏、阿毗达摩藏与小诵等先前所示的十五类，除去四部的其余佛语。

\begin{quoting}除去这长等的四部，\\其它的佛语，构成小部。\end{quoting}

如是依部为五种。

\subsection*{九分}

如何\textbf{依分为九种}？这一切以契经、应颂、记说、偈颂、自说、如是语、本生、未曾有法、毗陀罗等成为九类。这里，

\begin{enumerate}[nosep]
    \item 二部分别、义释、犍度、附随，经集中的吉祥经、宝经、那罗迦经、迅速经，与其它名为经的如来语，当知为\textbf{契经}。
    \item 一切有偈的经，当知为\textbf{应颂}。特别地，相应部中的整个有偈品亦然。
    \item 整个阿毗达摩藏，不含偈颂的经，与其它八支所未摄的佛语，当知为\textbf{记说}。
    \item 法句、长老偈、长老尼偈、经集中的不名为经的纯偈颂，当知为\textbf{偈颂}。
    \item 与以喜悦之智所造的偈颂相关的八十二经\footnote{现为八十经。}，当知为\textbf{自说}。
    \item 以「世尊如是说」等方法转起的一百一十经\footnote{现为一百一十二经。}，当知为\textbf{如是语}。
    \item 确然本生等五百五十本生\footnote{现为五百四十七本生。}，当知为\textbf{本生}。
    \item 以「诸比丘！于阿难中，有这四希有、未曾有法」{\bluefootnotesize 〔长部第 2.209 段〕}等方法转起一切希有、未曾有法相关的经，当知为\textbf{未曾有法}。
    \item 小毗陀罗、大毗陀罗、正见、帝释所问、行分析\footnote{行分析 \textit{Saṅkhāra-bhājaniya}：暂未找到此经，待考。}、大满月经等一切，在获得知识和满足后一再发问的经，当知为\textbf{毗陀罗}。
\end{enumerate}

如是依分为九种。

\subsection*{八万四千法蕴}

如何\textbf{依法蕴为八万四千种}？且这一切佛语，

\begin{quoting}我从佛陀取了八万二千，从比丘二千，\\由我转起之法，计八万四千。{\bluefootnotesize 〔长老偈第 1027 颂〕}\end{quoting}

如是，以所阐明的法蕴，成八万四千类。这里，一节的经为一法蕴。多节者，以节为法蕴之数。在偈颂的组合中，问题之问为一法蕴，答为一。在阿毗达摩中，一一三法、二法的分析与一一心番的分析，为一一法蕴。在律中，有事由、有总纲、有词的分析、有间犯、有犯、有无犯、有三法断\footnote{三法断 \textit{tikaccheda}：待考。}，这里一一部分当知为一一法蕴。如是依法蕴为八万四千种。

如是，这佛语由不分类，依味为一种，由分类，依法律等为二种等类，被以大迦叶为首的自在之众所结集——这是法、这是律，这是最初的佛语、这是中间的佛语、这是最后的佛语，这是律藏、这是经藏、这是阿毗达摩藏，这是长部……小部，这些是契经等九分，这些是八万四千法蕴——确定了这品类而结集。不仅如此，还确定了其它在三藏中出现的总颂之摄、品之摄、省略之摄、一集二集等的集之摄、相应之摄、五十之摄等多种摄的品类，经七月完成结集。

而在结集的终了，以「大迦叶长老所造的这十力的教法，堪能转起五千年的时量」而生起欢喜，似有掌声鼓起，大地以水为边界而震动、完全震动、完全遍震、完全遍动，且有许多希有现起。此即第一次大结集，它在世间：

\begin{quoting}由五百所造，因此为「五百」，\\且由长老所造故，而被称为「长老者」。\end{quoting}
